\section{Diferenciación de sucesiones de funciones}

\begin{theorem}
  Sea ${(f_n)}_{n \in \N}$ una sucesión de funciones derivables en $[a, b]$. Si existe $c \in [a, b]$ tal que ${(f_n(c))}_{n \in \N}$ converge y las derivadas $f_n'$ convergen uniformemente en $[a, b]$ a una función $g$, entonces $f_n \rightrightarrows f$ en $[a, b]$ para alguna función $f$, y $f' = g$. Es decir,
  \[
    {\left( \lim_{n \to +\infty} f_n \right)}' = \lim_{n \to +\infty} f_n'
  \]
  siempre que las derivadas converjan uniformemente (y las $f_n$ converjan en algún punto).
\end{theorem}

\begin{proof}
  \emph{Paso 1: ${(f_n)}$ converge uniformemente.} Aplicando el teorema del valor medio a $f_m - f_n$ en el intervalo de extremos $x$ y $c$, existe $d$ entre $x$ y $c$ tal que
  \[
    \bigl( f_m(x) - f_n(x) \bigr) - \bigl( f_m(c) - f_n(c) \bigr) = (x - c) \bigl( f_m'(d) - f_n'(d) \bigr).
  \]
  Luego, para todo $x \in [a, b]$,
  \[
    |f_m(x) - f_n(x)| \leq |f_m(c) - f_n(c)| + (b - a) \sup_{[a, b]} |f_m' - f_n'|.
  \]
  Como ${(f_n(c))}$ converge y ${(f_n')}$ es uniformemente de Cauchy, dado $\e > 0$ ambos términos son menores que $\frac{\e}{2}$ si $n, m$ son grandes, y entonces ${(f_n)}$ es uniformemente de Cauchy. Por lo tanto $f_n \rightrightarrows f$ para alguna $f$.

  \emph{Paso 2: $f' = g$.} Fijemos $x_0 \in [a, b]$ y definamos, para $x \in [a, b] \setminus \{x_0\}$,
  \[
    q_n(x) = \frac{f_n(x) - f_n(x_0)}{x - x_0}, \qquad q(x) = \frac{f(x) - f(x_0)}{x - x_0}.
  \]
  Por el teorema del valor medio aplicado a $f_m - f_n$ entre $x_0$ y $x$, existe $d$ entre $x_0$ y $x$ con $q_m(x) - q_n(x) = f_m'(d) - f_n'(d)$, de modo que $|q_m(x) - q_n(x)| \leq \sup_{[a, b]} |f_m' - f_n'|$ para todo $x$. Entonces ${(q_n)}$ es uniformemente de Cauchy en $[a, b] \setminus \{x_0\}$, y como $q_n \to q$ puntualmente, $q_n \rightrightarrows q$. Además $\lim_{x \to x_0} q_n(x) = f_n'(x_0)$. Por el teorema de intercambio de límites,
  \[
    \lim_{x \to x_0} q(x) = \lim_{n \to \infty} f_n'(x_0) = g(x_0),
  \]
  es decir, $f'(x_0) = g(x_0)$.
\end{proof}

\begin{corollary}
  Sean $f_n$ derivables en $[a, b]$. Si $\sum f_n(c)$ converge para algún $c \in [a, b]$ y $\sum f_n' \rightrightarrows g$, entonces $\sum f_n \rightrightarrows f$ en $[a, b]$, con $f$ derivable y $f' = g$.
\end{corollary}

\begin{proof}
  Se aplica el teorema a las sumas parciales $S_n = f_1 + \cdots + f_n$.
\end{proof}

\section{Series de potencias}

Una serie de la forma $\sum_{n \geq 0} a_n {(x - x_0)}^n$ se llama \emph{serie de potencias}. Por simplicidad consideramos $\sum_{n \geq 0} a_n x^n$; el caso general se reduce a éste mediante el cambio de variable $x \mapsto x - x_0$.

\begin{theorem}[Cauchy--Hadamard]
  Sea $\sum_{n \geq 0} a_n x^n$ una serie de potencias y sea $r \in [0, +\infty]$ definido por $\frac{1}{r} = \limsup_{n \to \infty} \sqrt[n]{|a_n|}$ (con $\frac{1}{0} = +\infty$ y $\frac{1}{+\infty} = 0$). Entonces la serie converge absolutamente si $|x| < r$ y diverge si $|x| > r$. El número $r$ se llama \emph{radio de convergencia}.
\end{theorem}

\begin{proof}
  Sea $\rho = \limsup \sqrt[n]{|a_n|} = \frac{1}{r}$. Para $x \neq 0$, $\limsup \sqrt[n]{|a_n x^n|} = |x| \rho$.

  Si $|x| \rho < 1$, es decir $|x| < r$ (esto incluye $\rho = 0$, donde converge para todo $x$), por el criterio de la raíz la serie converge absolutamente. Si $|x| \rho > 1$, es decir $|x| > r$ (esto incluye $\rho = +\infty$ y $x \neq 0$), existen infinitos $n$ con $\sqrt[n]{|a_n x^n|} > 1$, de modo que $|a_n x^n| > 1$ para infinitos $n$, el término general no tiende a $0$ y la serie diverge.
\end{proof}

\begin{theorem}
  Una serie de potencias converge uniformemente en todo intervalo compacto contenido en su intervalo de convergencia.
\end{theorem}

\begin{proof}
  Sea $r > 0$ el radio de convergencia de $\sum a_n x^n$ y sea $0 < s < r$. Para $x \in [-s, s]$ vale $|a_n x^n| \leq |a_n| s^n$, y como $\sum |a_n| s^n$ converge (la serie converge absolutamente en $s$), por el criterio de Weierstrass $\sum a_n x^n$ converge uniformemente en $[-s, s]$. Todo intervalo compacto contenido en $(-r, r)$ está contenido en algún $[-s, s]$.
\end{proof}

\begin{corollary}
  La función $f\colon (-r, r) \to \R$, $f(x) = \sum a_n x^n$, es continua.
\end{corollary}

\begin{proof}
  Dado $x_0 \in (-r, r)$, existe $s$ con $|x_0| < s < r$. En $[-s, s]$ la serie es un límite uniforme de polinomios, que son continuos, y por lo tanto $f$ es continua en $x_0$.
\end{proof}

Una serie de potencias puede no converger uniformemente en todo su intervalo de convergencia $(-r, r)$.

\begin{eg}
  La serie $\sum x^n$ no converge uniformemente en $(-1, 1)$. En efecto, si $\sum f_n \rightrightarrows f$ entonces $f_{n+1} = S_{n+1} - S_n \rightrightarrows 0$, pero $\sup_{x \in (-1, 1)} |x^n| = 1$ para todo $n$.
\end{eg}

\begin{lemma}
  Sea $\sum_p \alpha_p$ una serie con sumas parciales acotadas: si $S_p = \alpha_1 + \cdots + \alpha_p$, existe $k > 0$ con $|S_p| \leq k$ para todo $p \in \N$. Si $b_1 \geq b_2 \geq \cdots \geq b_p \geq 0$, entonces
  \[
    |\alpha_1 b_1 + \cdots + \alpha_p b_p| \leq k b_1.
  \]
\end{lemma}

\begin{proof}
  Con $S_0 = 0$ y $\alpha_i = S_i - S_{i-1}$,
  \begin{align*}
    |\alpha_1 b_1 + \cdots + \alpha_p b_p| & = |S_1 b_1 + (S_2 - S_1) b_2 + \cdots + (S_p - S_{p-1}) b_p|              \\
                                           & = |S_1 (b_1 - b_2) + \cdots + S_{p-1} (b_{p-1} - b_p) + S_p b_p|          \\
                                           & \leq k \bigl( (b_1 - b_2) + \cdots + (b_{p-1} - b_p) + b_p \bigr) = k b_1. \qedhere
  \end{align*}
\end{proof}

\begin{theorem}[Abel]
  Sea $\sum a_n x^n$ con radio de convergencia $0 < r < +\infty$. Si $\sum a_n r^n$ converge, entonces $\sum a_n x^n$ converge uniformemente en $[0, r]$. En particular, $\lim_{x \to r^{-}} \sum a_n x^n = \sum a_n r^n$.
\end{theorem}

\begin{proof}
  Sea $\e > 0$. Como $\sum a_n r^n$ converge, existe $n_0$ tal que $|a_{n+1} r^{n+1} + \cdots + a_{n+p} r^{n+p}| < \e$ para todo $p \in \N$ y todo $n > n_0$. Fijemos $n > n_0$ y $x \in [0, r]$, y sean $\alpha_p = a_{n+p} r^{n+p}$, cuyas sumas parciales están acotadas en módulo por $k = \e$, y $b_p = {\left( \frac{x}{r} \right)}^{n+p}$, que cumplen $b_1 \geq b_2 \geq \cdots \geq 0$ pues $\frac{x}{r} \in [0, 1]$. Como $a_{n+p} x^{n+p} = \alpha_p b_p$, el lema da
  \[
    |a_{n+1} x^{n+1} + \cdots + a_{n+p} x^{n+p}| \leq \e \, b_1 = \e {\left( \frac{x}{r} \right)}^{n+1} \leq \e
  \]
  para todo $p$ y todo $x \in [0, r]$. Luego la serie converge uniformemente en $[0, r]$. Entonces su suma $f$ es continua en $[0, r]$, por ser límite uniforme de polinomios, y $\sum a_n r^n = f(r) = \lim_{x \to r^{-}} f(x)$.
\end{proof}

\begin{note}
  \begin{enumerate}[label=(\arabic*)]
    \item El mismo resultado vale en $-r$: si $\sum a_n {(-r)}^n$ converge, aplicando el teorema a la serie $\sum {(-1)}^n a_n x^n$ se obtiene que $\sum a_n x^n$ converge uniformemente en $[-r, 0]$.
    \item $\sum a_n x^n$ converge uniformemente en $(-r, r)$ si y sólo si converge en $r$ y en $-r$. En efecto, si converge en ambos extremos, por el teorema de Abel converge uniformemente en $[0, r]$ y en $[-r, 0]$, y por lo tanto en $[-r, r]$. Recíprocamente, si converge uniformemente en $(-r, r)$, para $\e > 0$ existe $n_0$ con $|a_{n+1} x^{n+1} + \cdots + a_{n+p} x^{n+p}| < \e$ para todo $x \in (-r, r)$, $n > n_0$ y todo $p$, y haciendo $x \to \pm r$ resulta que la serie es de Cauchy en $\pm r$.
  \end{enumerate}
\end{note}

\begin{theorem}
  Sea $[\alpha, \beta] \subset [-r, r]$ tal que $\sum a_n x^n$ converge en $\alpha$ y en $\beta$. Entonces
  \[
    \int_{\alpha}^{\beta} \sum_{n \geq 0} a_n x^n \,\dd x = \sum_{n \geq 0} \frac{a_n}{n + 1} \left( \beta^{n+1} - \alpha^{n+1} \right).
  \]
\end{theorem}

\begin{proof}
  Por el criterio de Weierstrass y el teorema de Abel (según si $\alpha$ o $\beta$ son $\pm r$ o están en $(-r, r)$), la serie converge uniformemente en $[\alpha, \beta]$. Por el teorema de integración de series de funciones, puede integrarse término a término.
\end{proof}

Si $\sum a_n x^n$ no converge en un extremo $r$ de su intervalo de convergencia, igualmente puede integrarse término a término la integral impropia $\int_0^r \sum a_n x^n \,\dd x$, siempre que $\sum \frac{a_n r^{n+1}}{n+1}$ converja. En efecto, la serie $\sum \frac{a_n}{n+1} t^{n+1}$ tiene el mismo radio $r$ y converge en $r$, de modo que por el teorema de Abel,
\[
  \int_0^r \sum a_n x^n \,\dd x = \lim_{t \to r^{-}} \int_0^t \sum a_n x^n \,\dd x = \lim_{t \to r^{-}} \sum \frac{a_n}{n + 1} t^{n+1} = \sum \frac{a_n r^{n+1}}{n + 1}.
\]

\begin{theorem}
  La función $f(x) = \sum a_n x^n$ es derivable en todo el intervalo de convergencia $(-r, r)$, y $f'(x) = \sum n a_n x^{n-1}$, serie que tiene el mismo radio de convergencia.
\end{theorem}

\begin{proof}
  Como $\sqrt[n]{n} \to 1$, se tiene $\limsup \sqrt[n]{n |a_n|} = \limsup \sqrt[n]{|a_n|}$. Luego la serie $\sum n a_n x^n$, y por lo tanto $\sum n a_n x^{n-1}$ (que para $x \neq 0$ es $\frac{1}{x} \sum n a_n x^n$), tiene el mismo radio de convergencia $r$.

  Sea $0 < s < r$. La serie de derivadas $\sum n a_n x^{n-1}$ converge uniformemente en $[-s, s]$ por el teorema de convergencia uniforme de series de potencias, y $\sum a_n x^n$ converge en $c = 0$. Por el corolario del teorema de diferenciación de series de funciones, $f$ es derivable en $[-s, s]$ con $f'(x) = \sum n a_n x^{n-1}$. Como $s < r$ es arbitrario, vale en todo $(-r, r)$.
\end{proof}

\begin{corollary}
  Si $f(x) = \sum a_n x^n$ tiene radio de convergencia $r > 0$, entonces $f$ tiene derivadas de todos los órdenes en $(-r, r)$, que se calculan derivando término a término:
  \[
    f^{(k)}(x) = \sum_{n \geq k} n (n - 1) \cdots (n - k + 1)\, a_n x^{n-k}.
  \]
  En particular, $a_k = \dfrac{f^{(k)}(0)}{k!}$, y la serie de potencias que converge a $f$ en $(-r, r)$ es su serie de Taylor alrededor de $0$.
\end{corollary}

\section{Familias equicontinuas}

Dada una familia $E$ de funciones continuas con dominio común, queremos saber cuándo podemos garantizar que toda sucesión ${(f_n)}_{n \in \N} \subset E$ tiene una subsucesión uniformemente convergente.

A diferencia de lo que ocurre con sucesiones numéricas, no alcanza con que la familia esté uniformemente acotada. Por ejemplo, sea $f_n\colon [0, 1] \to \R$, $f_n(x) = x^n (1 - x^n)$. Entonces $0 \leq f_n \leq \frac{1}{4}$ y $f_n \to 0$ puntualmente, pero $f_n \bigl( 2^{-1/n} \bigr) = \frac{1}{2} \cdot \frac{1}{2} = \frac{1}{4}$ para todo $n$. Luego ninguna subsucesión puede converger uniformemente a $0$ y, como el límite uniforme debería coincidir con el puntual, ninguna subsucesión converge uniformemente.

\begin{definition}[Equicontinuidad]
  Sean $X \subset \R$ y $E$ una familia de funciones $f\colon X \to \R$. Dado $x_0 \in X$, decimos que $E$ es \emph{equicontinua en $x_0$} si para todo $\e > 0$ existe $\delta > 0$ tal que $|x - x_0| < \delta \Rightarrow |f(x) - f(x_0)| < \e$ para toda $f \in E$ (el mismo $\delta$ para todas las funciones). Decimos que $E$ es \emph{equicontinua} si lo es en todo punto de $X$.
\end{definition}

\begin{theorem}
  Sean $f_n\colon X \to \R$ continuas y $f_n \rightrightarrows f$. Entonces la sucesión ${(f_n)}_{n \in \N}$ es equicontinua.
\end{theorem}

\begin{proof}
  La función $f$ es continua por ser límite uniforme de continuas. Sean $x_0 \in X$ y $\e > 0$. Existe $n_0$ tal que $|f_n(x) - f(x)| < \frac{\e}{3}$ para todo $n > n_0$ y todo $x \in X$. Por la continuidad de $f$ existe $\delta_0 > 0$ tal que $|x - x_0| < \delta_0 \Rightarrow |f(x) - f(x_0)| < \frac{\e}{3}$, y por la de cada $f_i$ existe $\delta_i > 0$ tal que $|x - x_0| < \delta_i \Rightarrow |f_i(x) - f_i(x_0)| < \e$, para $i \in \{1, \dots, n_0\}$. Sea $\delta = \min\{ \delta_0, \delta_1, \dots, \delta_{n_0} \}$. Si $|x - x_0| < \delta$ y $n \leq n_0$, entonces $|f_n(x) - f_n(x_0)| < \e$. Si $n > n_0$,
  \[
    |f_n(x) - f_n(x_0)| \leq |f_n(x) - f(x)| + |f(x) - f(x_0)| + |f(x_0) - f_n(x_0)| < \e. \qedhere
  \]
\end{proof}

\begin{eg}
  La sucesión $f_n(x) = \dfrac{\sin(nx)}{n}$ es equicontinua en $\R$, pues $|f_n(x)| \leq \frac{1}{n}$ implica $f_n \rightrightarrows 0$.
\end{eg}

\begin{eg}
  La sucesión $f_n(x) = nx$ no es equicontinua en ningún $x_0 \in \R$. Sea $\e = \frac{1}{2}$ y $\delta > 0$. Existe $n \in \N$ con $\frac{1}{n} < \delta$, y el punto $x = x_0 + \frac{1}{n}$ cumple $|x - x_0| = \frac{1}{n} < \delta$, pero $|f_n(x) - f_n(x_0)| = n \cdot \frac{1}{n} = 1 \geq \e$.
\end{eg}
